\section{Full results}
\label{app:full}

Each cell is the median over runs. Throughput is in operations per second (k $=10^3$, M $=10^6$, G $=10^9$). Ratios above 1 favour HStore. ``Agree'' is \emph{yes} when both systems returned the same checksum in every run.

\subsection{Scale 1, asynchronous commits}

\begin{center}\footnotesize\input{gen/main-s1-async}\end{center}

\begin{center}\footnotesize Latency (median of each percentile over runs)\\[2pt]\input{gen/latency-s1-async}\end{center}

\begin{center}\footnotesize\input{gen/writes-s1-async}\end{center}

\begin{center}\footnotesize\input{gen/memory-s1-async}\end{center}

\subsection{Scale 1, synchronous commits}

\begin{center}\footnotesize\input{gen/main-s1-sync}\end{center}

\begin{center}\footnotesize Latency (median of each percentile over runs)\\[2pt]\input{gen/latency-s1-sync}\end{center}

\begin{center}\footnotesize\input{gen/writes-s1-sync}\end{center}

\begin{center}\footnotesize\input{gen/memory-s1-sync}\end{center}

\subsection{Scale 4, asynchronous commits}

\begin{center}\footnotesize\input{gen/main-s4-async}\end{center}

\begin{center}\footnotesize Latency (median of each percentile over runs)\\[2pt]\input{gen/latency-s4-async}\end{center}

\begin{center}\footnotesize\input{gen/writes-s4-async}\end{center}

\begin{center}\footnotesize\input{gen/memory-s4-async}\end{center}

\clearpage
\section{Run-to-run variability}
\label{app:variability}

The tables below give the minimum and maximum over runs and the coefficient of variation (standard deviation over mean) of every throughput workload, followed by the value of every run.

\subsection{Scale 1, asynchronous commits}
\begin{center}\footnotesize\input{gen/ranges-s1-async}\end{center}
{\footnotesize\input{gen/runs-s1-async}}

\subsection{Scale 1, synchronous commits}
\begin{center}\footnotesize\input{gen/ranges-s1-sync}\end{center}
{\footnotesize\input{gen/runs-s1-sync}}

\subsection{Scale 4, asynchronous commits}
\begin{center}\footnotesize\input{gen/ranges-s4-async}\end{center}
{\footnotesize\input{gen/runs-s4-async}}

\clearpage
\section{Reproducing the results}
\label{app:repro}

The harness is the \texttt{benchmarks} module of the HStore repository. HyperGraphDB and JE are built from source at the commits named in \cref{tab:config}.

\begin{verbatim}
benchmarks/install-hypergraphdb.sh      # builds HyperGraphDB 99485a1 and JE 5.0.73
SCALE=1 DURABILITY=async CACHE_MB=64 RUNS=5 THREADS=8 benchmarks/run.sh
SCALE=1 DURABILITY=sync  CACHE_MB=64 RUNS=5 THREADS=8 benchmarks/run.sh
SCALE=4 DURABILITY=async CACHE_MB=64 RUNS=3 THREADS=8 benchmarks/run.sh
\end{verbatim}

\texttt{run.sh} builds the module, then runs each system \texttt{RUNS} times in alternation, one JVM per run, with these flags:

\begin{verbatim}
java -Xms4g -Xmx4g -XX:+UseParallelGC \
     --add-opens java.base/java.lang=ALL-UNNAMED \
     --add-opens java.base/java.lang.ref=ALL-UNNAMED \
     --add-opens java.base/java.util=ALL-UNNAMED \
     io.hstore.bench.Comparison run --store <hstore|hypergraphdb> \
     --scale <s> --durability <async|sync> --threads 8 --history 64 --cache-mb 64
\end{verbatim}

Each run writes one JSON file with every measurement in this report. \texttt{Comparison report} turns a set of them into Markdown tables. The tables, figures and quoted numbers in this report were generated from the same files by \texttt{benchmarks/report/gen.py}, without manual transcription. The dataset statistics in \cref{tab:dataset} are printed by \texttt{io.hstore.bench.DatasetStats}, run with the same classpath. The raw files are in \texttt{benchmarks/report/results}, and the report is rebuilt with:

\begin{verbatim}
cd benchmarks/report
python3 gen.py results gen
tectonic report.tex
\end{verbatim}

\section{HStore settings}
\label{app:settings}

\begin{center}\footnotesize
\begin{tabular}{@{}llr@{}}
\toprule
Setting & Meaning & Value \\
\midrule
\texttt{page\_size} & tree node page size & 16\,KiB \\
segment size & pages per segment file & 16,384 (256\,MiB) \\
\texttt{node\_cache\_mb} & decoded node cache & 64\,MiB (set) \\
\texttt{write\_buffer\_mb} & cap on unwritten tree nodes & 64\,MiB \\
\texttt{checkpoint\_wal\_mb} & log or page bytes that trigger a checkpoint & 256\,MiB \\
\texttt{compaction\_live\_ratio} & compact segments below this live fraction & 0.5 \\
\texttt{history\_limit} & generations kept for time travel & 64 \\
\texttt{wal\_mode} & log page references, not page images & references \\
maintenance interval & background checkpoint and compaction pass & 500\,ms \\
\bottomrule
\end{tabular}
\end{center}
